\section{Przebieg eksperymentu}

\paragraph{Etap 0.} Konstrukcja bazowa (dwa odwzorowania trójkątne z interpolacji Lagrange'a)
dała automorfizm stopnia $16=(m-1)^2$, przyjęty przez weryfikator.

\paragraph{Iteracja 01.} Planista zauważył, że brak niestałego funkcjonału afinicznego zgodnego
na $p_i$ i $q_i$ wyklucza słowa krótkie, i zaplanował trzy zadania.
\begin{itemize}
\item zad1: przypadek $(2,2)$ wykluczony obliczeniowo (metoda I); przy okazji znaleziono
automorfizm stopnia 6, multistopień $(3,2)$.
\item zad2: dwa automorfizmy stopnia 6 o multistopniach $(2,3)$ i $(3,2)$ (ten drugi daje
to samo $F$ co rozwiązanie zad1, zob. tabela wyników -- identyczne wielomiany).
\item zad3: dowód lematu~\ref{lem:postac} i wykluczenia słów długości $0$ i $1$; wniosek
$\deg F\ge4$ i $\deg F$ złożony.
\end{itemize}

\paragraph{Iteracja 02.} Celem była niezależna kontrola wyniku negatywnego i spisanie dowodu.
\begin{itemize}
\item zad1: metoda II (pełna baza Gröbnera) potwierdziła brak rozwiązań wymiernych dla $(2,2)$;
test kontrolny nad ciałami liczbowymi wypadł pomyślnie.
\item zad2: spisanie i recenzja całego łańcucha $\deg F\ge6$; nie znaleziono luki merytorycznej,
doprecyzowano wchłanianie części afinicznych przez $\sigma$ i normalizacje.
\item zad3: czytelniejsze rozwiązanie stopnia 6 (32 wyrazy) oraz trzecie niezależne sprawdzenie
braku zgodnego funkcjonału. W siatce kierunków $|\cdot|\le3$ rozwiązanie istnieje tylko dla 6 par
$(a,e)$ (wynik obliczeniowy).
\end{itemize}

\paragraph{Iteracja 03.} Planista uznał pytanie za rozstrzygnięte i nie zlecił nowych zadań.
